\documentclass[12pt]{article}
\usepackage[scaled=0.75]{DejaVuSansMono}
\usepackage{fullpage}
\usepackage{amsfonts}
\usepackage{amsmath}
\usepackage[T1]{fontenc}
\usepackage[utf8]{inputenc}
\usepackage{enumitem}
\usepackage{graphicx}
%\usepackage[cache=false]{minted}
\usepackage{listings}
\usepackage{hyperref}
\hypersetup{
    colorlinks=true,
    linkcolor=blue,
    filecolor=magenta,      
    urlcolor=cyan,
}
\usepackage{color}
\definecolor{keywordcolor}{rgb}{0.7, 0.1, 0.1}   % red
\definecolor{tacticcolor}{rgb}{0.0, 0.1, 0.6}    % blue
\definecolor{commentcolor}{rgb}{0.4, 0.4, 0.4}   % grey
\definecolor{symbolcolor}{rgb}{0.0, 0.1, 0.6}    % blue
\definecolor{sortcolor}{rgb}{0.1, 0.5, 0.1}      % green
\definecolor{attributecolor}{rgb}{0.7, 0.1, 0.1} % red

\def\lstlanguagefiles{lstlean.tex}
%\lstset{language=lean,basicstyle=\ttfamily,breaklines=true}

\newtheorem{definition}{Definition}
\newtheorem{theorem}{Theorem}
\newtheorem{proposition}[theorem]{Proposition}
\newtheorem{lemma}[theorem]{Lemma}
\newtheorem{corollary}[theorem]{Corollary}

\title{MS-EV0029 Project Contribution Report}
\author{Risto Raveala}
\date{\today}

\begin{document}

\maketitle


\section*{Project overview}
As a part of the course MS-EV0029, I built a project that contains a proof of the Central Limit Theorem, using only results and definitions that are available in Mathlib. The Github repo is behind the following link:
\begin{center}
\textbf{Central Limit Theorem} \\
Github repository: \url{https://github.com/ristoraveala/course_project}
\end{center}

The goal of the project is to show, that if $X_1$, $X_2,\ldots$ is a sequence of independent identically distributed random variables such that $\mathbb{E} X_1^2 < \infty$, and such that the sequence is not a sequence of deterministic constants (i.e. $\textrm{Var} (X_1) \ne 0$), then
\[
\frac{1}{\sqrt{n \cdot \textrm{Var} (X_1)}} \sum_{k=1}^n (X_k - \mathbb{E} X_k) \overset{d}{\longrightarrow} \mathcal{N}(0,1), \qquad n \to \infty,
\]
where $\mathcal{N}(0,1)$ denotes the standard normal distribution. More precisely, the main theorem is of the project is:

\begin{theorem}[Central limit theorem]
Suppose that $(X_1,X_2,\ldots)$ is a sequence of iid random variables such that $\mathbb{E} X_1^2 < \infty$. Denote
\[
\mu_n := \textrm{Law} \left[\frac{1}{\sqrt{n \cdot \textrm{Var} (X_1)}} \sum_{k=1}^n (X_k - \mathbb{E} X_k) \right].
\]
Then in the topology of weak convergence of probability measures of $\mathbb{R}$, the sequence $(\mu_n)_{n\in\mathbb{N}}$ converges to the standard normal distribution.
\end{theorem}


\section*{Main contributions}
Below is a list of my most significant contributions made during the project:

\begin{enumerate}[label=\arabic*:, leftmargin=1cm]

    \item commit \href{https://github.com/ristoraveala/course_project/commit/880711fffb6b6ed5387d25e7a5b5876e41b2cb31}{880711f} \\
    \textbf{Description:}
    This contains the proof of most of the mathematical content one of the main lemmas, namely \lstinline{charFun_convergence}. The proof of \lstinline{charFun_convergence} is spread out in multiple commits, with the final commit dedicated to this proof being \href{https://github.com/ristoraveala/course_project/commit/596838fccb7eb732d234e8d88bad2c0218fa968c}{596838f}. The proof of this lemma was the most difficult task of the project, and it is difficult to identify one single commit in which the proof was finalized.

    
    \item commit \href{https://github.com/ristoraveala/course_project/commit/fc7fb82e8230f713eea08c8b9e075fd9af75da4c}{fc7fb82} \\ 
    Proved \lstinline{iid_sum_taylor_product}, from which the second key lemma, \lstinline{iid_sum_to_cexp} follows.

    \item commit \href{https://github.com/ristoraveala/course_project/commit/b6b62af2eea7d5f06709dfc11b21cb1c5c2c5514}{b6b62af} \\ 
    Proved the the second key lemma, \lstinline{iid_sum_to_cexp} which shows that the characteristic functions of the $\mu_n$ as defined above converge pointwise to the characteristic function of standard normal distribution.
    

    %Add more commits as needed
    %\item commit \href{https://the.url.of.fi/your/commit/or/pr}{abc123} \\ 
    %Description.
    %\lstinline{leanDeclarationName}
    %\emph{(blueprint-label)}
    
\end{enumerate}

The two main lemmas of the project were \lstinline{iid_sum_to_cexp} and \lstinline{charFun_convergence}. The second lemma required by far the most work in this project.

The statement of the lemma is that if the characteristic functions of a sequence of probability measures $(\nu_n)$ (on $\mathbb{R}$) converge pointwise to a characteristic function of a probability measure $\nu$ (on $\mathbb{R}$), then $\nu_n \to \nu$ as $n\to\infty$ in the topology of weak convergence of measures. The proof has many stages: at first we prove the statement for measures with densities with respect to the Lebesgue measure on $\mathbb{R}$. This is lemma \lstinline{charFun_convergence_with_densities}. Then we prove that the convolution of the measures with a gaussian measure has a density and does satisfy the statement of the lemma \lstinline{charFun_convergence_with_densities}. We go on to prove a uniform estimate \lstinline{lipschitz_bdd_unif_approx} which shows that for fixed bounded and Lipschitz continuous function $f \colon \mathbb{R} \to \mathbb{R}$, the convergence $\int f \, \textrm{d}\nu^\delta \to \int f \, \textrm{d}\nu$, where $\nu^\delta$ denotes the convolution of $\nu$ with the law of a centered Gaussian with variance $\delta$, converges uniformly with respect to the measure $\nu$. From these ingredients the lemma \lstinline{charFun_convergence} follows.

The proof of lemma \lstinline{iid_sum_to_cexp} is mathematically more straightforward, but requires some work in lean. The idea is that we show, as a consequence of multiple lemmas that build on this, that the characteristic function $\varphi_n$ of $(\sqrt{n \cdot \textrm{Var} (X_1)})^{-1} \sum_{k=1}^n (X_k - \mathbb{E} X_k)$ satisfies
\[
\varphi_n(t) = \left(1 - \frac{t^2}{2n} + o(1/n) \right)^n,
\]
which converges to $e^{-\frac{t^2}{2}}$, the characteristic function of a standard normal distribution at $t$, as $n\to\infty$.

\section*{Notes}

\subsubsection*{Collaboration:}
This project was made entirely by me.
\end{document}
